% =====================================================================
%  第九章  设计评价与不足
% =====================================================================
\chapter{设计评价与不足}

\section{已完成目标}
本次课程设计围绕 UNIX V6++ 内存管理的离散化展开，已完成课设提出的全部三大改造目标：

\begin{enumerate}
  \item \textbf{页表系统改造}。实现了每进程独立的页目录，进程创建时动态分配页目录并
        写入 \texttt{CR3}；每个进程保留一张私有 1\# 用户页表，核心页表与 0\# 用户页表
        全局共享；进程切换通过 \texttt{SwtchUStruct} 更新内核页表 user 项并刷新
        \texttt{CR3} 完成（第三章）。
  \item \textbf{物理页框分配改为位示图}。以 \texttt{BitMap} 替代连续空闲分区表，
        \texttt{BitMapAllocator} 通过位运算实现离散分配与释放；页表区与用户区分别由
        \texttt{KernelPageManager}/\texttt{UserPageManager} 管理；引入 \texttt{Page[]}
        页引用计数支撑共享页生命周期（第四章）。
  \item \textbf{离散化进程映像管理}。实现了逻辑页离散存放、fork 写时复制、exec 直接
        挂用户页表并构造参数栈、正文段共享（\texttt{Text::x\_caddr[]} 存物理页框号）
        四项要求（第五章）。
\end{enumerate}

在工程层面，本设计同时对构建与运行环境进行了适配，使改造后的内核能够完整构建、
引导启动并运行用户程序，从而具备了在本地验证各机制的条件。

\section{当前实现的优点}
归纳起来，当前实现有如下几个优点。

\textbf{第一，改造层次清晰。}整个设计自然地分为页表层、页框分配层、进程生命周期层
三个相互独立又彼此耦合的层次。页表改造、位示图分配与引用计数各自负责明确的职责：
页表负责虚\,$\to$\,物映射，位示图负责页框的分配回收，引用计数负责描述共享物理页的
生命周期，COW 则把 fork 后的共享页面延迟复制到真正写入时再执行。这种拆分使得系统
可以沿"页表、页框、进程生命周期"三层结构加以解释，而不必把所有逻辑混在某一个函数里。

\textbf{第二，保留了用户态连续地址空间的抽象。}对用户程序而言，text、data、heap 和
stack 仍按熟悉的方式组织，无需感知物理页框是否连续；对内核而言，真实物理内存可以按
4KB 页框离散分配，并由页表重新拼接成连续虚拟空间。这正是分页机制的核心价值，也符合
课设"离散化存储"的目标。

\textbf{第三，fork 采用写时复制。}相比简单复制整个用户映像，COW 既减少了 fork 时的
物理内存复制开销，也能自然地处理父子进程共享页面、后续写入分离、text 段只读共享等
情形；引用计数的加入使页面释放不再只看某一个进程的页表项，而是能判断该物理页是否仍
被其他进程使用。

\textbf{第四，工程可运行性较强。}本设计处理并记录了工具链适配、Makefile 改写、
bootloader 加载、镜像布局等问题，使系统不致停留在代码层面，而是能够启动到 shell 并
运行用户程序。对操作系统课程设计而言，"能构建、能启动、能验证"本身就是实现质量的
重要组成部分。

\section{不足与后续改进}
当前实现仍存在若干不足，在报告中明确说明如下。

\textbf{其一，swap 相关遗留路径未完全清理。}核心用户页框分配、fork/COW、exec/exit 等
运行时主路径已经转向离散页框模型，不再依赖旧的连续进程映像换入换出作为主机制；但原
框架中仍保留 \texttt{SwapperManager} 的若干遗留路径，例如 text 副本、僵尸进程 u 区
保存以及内存不足时的旧接口。这些遗留路径在当前验证范围内并非主路径，但从系统设计
完整性看，后续仍应进一步梳理和清理。

\textbf{其二，ATA 磁盘 I/O 采用同步 PIO 替代 DMA。}这一选择提高了本地 QEMU 验证的
可靠性——因为在离散页表模型下，缓冲区的线性地址不能再通过简单掩码直接得到 DMA 所需
的物理地址；但 PIO 的效率低于 DMA，也没有真正解决"离散页表模型下如何为 DMA 提供正确
物理缓冲区描述"的问题。后续若继续完善，可基于页表查询或专门的 DMA 缓冲区机制，为
ATA DMA 构造可靠的物理地址或 PRD 表。

\textbf{其三，内存不足处理较为粗糙。}bitmap 分配器能够判断页框是否可用，但很多上层
路径仍倾向于假设分配成功，或在失败时直接 panic。一个更完整的系统应当在用户页框不足
时向系统调用返回错误，或实现更明确的进程回收、页面回收、失败回滚策略。由于课设明确
不使用盘交换区，内存不足时如何优雅失败也应成为后续改进重点。

\textbf{其四，自动化测试与文件系统构建工具的跨平台支持不足。}当前的验证以手动用例
为主，缺乏回归测试套件；同时，V6++ 文件系统镜像的构建仍依赖 Windows 专有工具，本地
跨平台生成镜像尚需额外工作。后续可补充自动化测试，并实现基于 \texttt{mtools} 或自研
写入器的跨平台镜像构建方案。
